\section{Introduction}\label{sec:intro}

A contractible graph can be transformed into a single vertex by adding
and deleting vertices and edges whose corresponding neighbourhoods are
contractible. Requiring a reduction by deletions alone, and imposing the
same requirement recursively on the neighbourhoods, gives the class of
strong I-contractible graphs. What is the least order of an I-contractible graph that is not strong
I-contractible? Does replacing contractibility by acyclicity over
$\mathbb F_2$ lower this minimum?

Ivashchenko introduced contractible transformations of graphs
\cite{Ivashchenko1993,Ivashchenko1994a}. The clique complex $\cl(G)$ has
the cliques of $G$ as its simplices. The class $\I$ of I-contractible
graphs is precisely the class of graphs with contractible clique complex
\cite[Corollary~1.5]{DochtermannMatsushita2025}.
Strong I-contractible graphs, denoted by $\IS$, have collapsible clique
complexes \cite[Theorem~1]{EFH2022}; hence $\IS\subseteq\I$.

Dochtermann, Espinoza, Fr\'ias-Armenta and Hern\'andez
\cite{DEFH2023} counted several contractible and dismantlable families
through order $11$. They also exhibited a $15$-vertex graph $\Dunce$ whose
clique complex is a flag triangulation of the dunce hat, with
$\Dunce\in\I\setminus\IS$, and a one-edge extension $\Dunce^+\in\IS$.
In our notation, their conjectures assert that these examples have
minimum order:

\begin{quote}
\textbf{Conjecture 14 of \cite[p.~8]{DEFH2023}.}
The graph $\Dunce$ is smallest by number of vertices among graphs with
trivial reduced $\mathbb F_2$-homology that are not strong
I-contractible. In particular, it is smallest among the graphs in
$\I\setminus\IS$.

\smallskip
\textbf{Conjecture 15 of \cite[p.~9]{DEFH2023}.}
The graph $\Dunce^+$ is smallest by number of vertices among graphs
$G\in\IS$ having an edge $uv$ for which the common neighbourhood
$N_G(u,v)$ belongs to $\IS$ but $G-uv\notin\IS$.
\end{quote}

We refute both conjectures. Three graphs on $13$ vertices have clique
complexes homotopy equivalent to the Moore space $M(\ZZ/3,1)$.
Their reduced homology vanishes over $\mathbb F_2$, whereas their first
integral homology is $\ZZ/3$. They therefore lie outside $\I$ and
refute the first assertion of Conjecture~14. A graph on $14$ vertices
belongs to $\I\setminus\IS$ and refutes its second assertion. Its
membership in $\I$ follows from one admissible edge addition followed
by a strong reduction. The same extension refutes Conjecture~15.
Section~\ref{sec:examples} gives the constructions and their proofs.

Our second contribution is an extension of the cardinality table,
together with algorithms and proofs that justify its computation.
We count connected graphs up to isomorphism in the contractible and
dismantlable families through order $13$, and in the connected
$K_4$-free universe at order $14$. Integral and $\mathbb F_2$
acyclicity are recorded separately. The integral and I-contractible
columns are derived from the $\mathbb F_2$ census and the classification
of its non-strong graphs. The three torsion examples are
exactly the $\mathbb F_2$-acyclic graphs outside $\IS$ through order
$13$. Consequently, the minimum order for the first assertion of
Conjecture~14 is $13$, while the minimum for its second assertion and
for the property in Conjecture~15 is $14$. At order $14$ there are
exactly eight $K_4$-free members of $\I\setminus\IS$; the fewest
edges among these eight is $44$.

The table also records the hierarchy of $k$-dismantlable graphs
introduced by Fieux and Jouve \cite{FieuxJouve2020}.
Throughout the listed universes, $2$-Dis and $3$-Dis coincide with
the strong vertex I-contractible class $\ISV$. The first separation
between $1$-Dis and $2$-Dis occurs at order $12$, with $20$ graphs.

The computation uses exact recursions for the strong classes and the
dismantling levels. At each stage these recursions consider all possible
first deletions. The remaining acyclic graphs are classified by explicit
contractible transformations and integral homology. For $K_4$-free
graphs, vertex neighbourhoods have clique complexes of dimension at most one and
edge neighbourhoods are discrete; this gives local admissibility criteria
and reduces the order-$14$ classification to graphs with no
admissible deletion.

The $13$-vertex construction also attains the minimum order for odd
torsion in the first integral homology of a flag complex, by comparison
with Adamaszek's classification \cite{Adamaszek2014}.
The three $14$-vertex graphs constructed below have contractible,
non-collapsible clique complexes. The computer-assisted classification through order
$13$ shows that $14$ is the minimum number of vertices of such a
flag complex (Corollary~\ref{cor:minimality}). Earlier examples arise
from Bing's house \cite{ChenYauYeh2001} and from flag triangulations
of the dunce hat \cite{FieuxJouve2020,DEFH2023}.
Our classification of order-$14$
graphs is restricted to the $K_4$-free class; the minimum by number
of vertices follows from the unrestricted census through order $13$.

\medskip\noindent\textbf{Organization.}
Section~\ref{sec:prelim} fixes the graph operations and homological
conventions. Section~\ref{sec:examples} proves the refutations using
explicit examples. Sections~\ref{sec:algorithms} and~\ref{sec:searches}
justify the enumeration and derive the table and the minimum orders.
The appendix supplies the finite mathematical data used in the
constructions.

\medskip\noindent\textbf{Conventions and notation.}
Graphs are finite and simple. Their order is their number of vertices.
Homology refers to the clique complex and is reduced unless stated
otherwise. The empty graph belongs to none of the contractible or
acyclic classes considered here. Cardinalities count isomorphism classes.
